%%% Преамбула ГОСТ 7.32
%%% Артемий Владин
%%% https://github.com/temavladin


%% Формат бумаги, размер шрифта (12pt или 14pt)
%% documentclass a4paper 14pt extarticle


%% Локализация
\usepackage[english,russian]{babel}
\usepackage{indentfirst} %% Красная строка
\addto\captionsrussian{\renewcommand{\figurename}{Рисунок}}
\addto\captionsrussian{\renewcommand{\tablename}{Таблица}}
\addto\captionsrussian{\renewcommand{\contentsname}{\hfill \normalsize \bfseries СОДЕРЖАНИЕ \hfill}}


%% Основной шрифт
\usepackage{fontspec}
\setmainfont[Ligatures=TeX]{Times New Roman}

%% Машинописный шрифт
\setmonofont{Courier New}
%\setmonofont{Consolas}[Scale=0.9]
%\setmonofont{Fira Mono}

%% Математический шрифт (Рекомендуемые: STIX или XITS)
\usepackage{mathtools}
\usepackage[warnings-off={mathtools-colon,mathtools-overbracket}]{unicode-math}
%\setmathfont{STIX Two Math}
%\setmathfont{STIX Two Math}[range={it/greek->up,bfit/greek->bfup}]


%% Поля документа
\usepackage[left=30mm, right=15mm, top=20mm, bottom=20mm]{geometry}
%% Нумерация страниц
\pagestyle{plain}

%% Абзацный отступ
\setlength{\parindent}{1.25 cm}
%% Межабзацный интервал
\setlength{\parskip}{0pt}

\usepackage{setspace}
%% Полуторный межстрочный интервал
\onehalfspacing


%% Настройка переноса слов
\usepackage{microtype}
\pretolerance=-1 %% Отключение попытки выравнивания строки без переносов
\tolerance=300 %% Позволяет избежать длинных пробелов
\emergencystretch=3em %% Аварийное растяжение межсловных пробелов
\exhyphenpenalty=150 %% Штраф для переносов после дефиса
\doublehyphendemerits=10000 %% Запрет на два переноса подряд в двух соседних строках
\finalhyphendemerits=10000 %% Запрет переноса слова в предпоследней строке абзаца

%% Отключение переноса слов
%\tolerance=1 %% Допуск на разряжённость
%\emergencystretch=\maxdimen %% Растяжение межсловных пробелов
%\hyphenpenalty=10000 %% Запрет на переносы слов
%\hbadness=10000 %% Мера разряжённости


%% Введение рисунков
\usepackage{graphicx}
%% Плавающие объекты
\usepackage{float}
%% Подключение цветов
\usepackage{xcolor}
%% Настройка подписей
\usepackage{subfig}


%% Подключение подписей к объектам
\DeclareCaptionLabelFormat{gost}{#1~#2} %% Формат метки: "Рисунок 1", "Таблица 2"
\DeclareCaptionLabelSeparator{gost}{~---~} %% Разделитель между меткой и текстом подписи
\usepackage{caption} %% Подписи к плавающим объектам

%% Рисунки по центру без переносов
\DeclareCaptionFormat{gost-figure}{%
\begingroup
\centering
\hyphenpenalty=10000
\exhyphenpenalty=10000
#1#2#3%
\endgroup
}

%% Таблицы слева без переносов
\DeclareCaptionFormat{gost-table}{%
\begingroup
\raggedright
\hyphenpenalty=10000
\exhyphenpenalty=10000
#1#2#3%
\endgroup
}

\captionsetup*[figure]{format=gost-figure, labelformat=gost, labelsep=gost}
\captionsetup*[table]{format=gost-table, labelformat=gost, labelsep=gost, singlelinecheck=false}


%% Стили заголовков
\usepackage{xparse} %% Создание своих команд
\NewDocumentCommand{\gostheading}{s m}{ %% Команда для заголовка по ГОСТ
\clearpage
\IfBooleanTF{#1}
{% Звёздочка ЕСТЬ: просто печатаем заголовок (без добавления в содержание)
\begin{center}
\bfseries\MakeUppercase{#2}
\end{center}
}
{% Звёздочки НЕТ: добавляем в оглавление и создаём якорь
\phantomsection
\addcontentsline{toc}{section}{#2}
\begin{center}
\bfseries\MakeUppercase{#2}
\end{center}
}
\par
}


%% Стили заголовков
\usepackage{titlesec}
\titleformat{\section}[hang]
{\normalsize \bfseries \hyphenpenalty=10000} %% Формат
{\thesection}{1 em} %% Нумерация и отступ
{} %% Любой доп. код

\titleformat{\subsection}[hang]
{\normalsize \bfseries \hyphenpenalty=10000}
{\thesubsection}{1 em}
{}
	
\titleformat{\subsubsection}[hang]
{\normalsize}
{\thesubsubsection}{1 em}
{}

%% Отступы от заголовка
\titlespacing*{\section}{12.5 mm}{0.5\baselineskip}{0.1\baselineskip}
\titlespacing*{\subsection}{12.5 mm}{0.5\baselineskip}{0.1\baselineskip}
\titlespacing*{\subsubsection}{12.5 mm}{\parskip}{-\parskip}


%% Настройки содержания
\usepackage{titletoc}

\titlecontents{section}
[2 em] % Общий отступ от левого края
{\vspace{0.3\baselineskip}} % Вертикальный промежуток 
{\contentslabel{2 em}} % Нумерованные
{\hspace*{-2 em}} % Ненумерованные
{\titlerule*[0.5pc]{.}\contentspage} % Точки

\titlecontents{subsection}
[4.9 em]
{}
{\contentslabel{2.7em}}
{\hspace*{2.7em}}
{\titlerule*[0.5pc]{.}\contentspage}

\titlecontents{subsubsection}
[8.6 em]
{}
{\contentslabel{3.5 em}}
{\hspace*{3.5 em}}
{\titlerule*[0.5pc]{.}\contentspage}


%% Стиль списков
\usepackage{enumitem}
\AddEnumerateCounter{\asbuk}{\@asbuk}{\cyrm}
%% Списки начинаются с красной строки (12.5 мм)
\setlist[itemize]{wide=\parindent, leftmargin=*, label={---}, itemsep=0pt, topsep=0pt, parsep=0pt, partopsep=0pt}
\setlist[enumerate]{wide=\parindent, leftmargin=*, itemsep=0pt, topsep=0pt, parsep=0pt, partopsep=0pt}
\setlist[enumerate,1]{label=\arabic*.}
\setlist[enumerate,2]{label=\asbuk*)}


%% Гиперссылки
\usepackage{hyperref}
\hypersetup{
unicode,
colorlinks,
citecolor=black,
filecolor=black,
linkcolor=black,
urlcolor=blue %% Ссылки в интернет синие
}


%% Настройка таблиц
\renewcommand{\arraystretch}{1.3} %% Высота строк таблицы
%\setlength{\tabcolsep}{1 em} % Отступы слева и справа от содержимого ячейки


%% Авто датировка
\usepackage[ddmmyy]{datetime} %% Дата \today
\renewcommand{\dateseparator}{.} %% Разделитель точка


%% Дополнительные пакеты
\usepackage{amsmath, amsthm, amsfonts} %% Математика формулы
%\usepackage{amssymb} %% Этот пакет не работает вместе с unicode-math
\usepackage{icomma} %% Умная запятая для чисел
\usepackage{textcomp}
\usepackage{tabularray}
\usepackage{rotating}
\usepackage{tabularx}
\usepackage{array}
\usepackage{ragged2e}
\usepackage{makecell} %% Переносы в клетках таблицы
\usepackage{hhline} %% Улучшенные горизонтальные линии в таблицах
\usepackage{multirow} %% Ячейки в несколько строчек в таблицах
\usepackage{longtable} %% Многостраничные таблицы


%% Листинг
\usepackage{fvextra} %% Улучшенный Verbatim


%% Настройка библиографии
\makeatletter

% 1. Переопределяем окружение thebibliography, чтобы убрать заголовок
\renewenvironment{thebibliography}[1]{
	\par % начинаем с нового абзаца
	\setcounter{mybibcntr}{0} % сбрасываем счетчик источников
}{
	\par % заканчиваем абзацем
}

% 2. Переопределяем команду \bibitem для вывода "Цифра. Текст" как абзац
\newcounter{mybibcntr}
\renewcommand{\bibitem}[1]{%
	\stepcounter{mybibcntr}%
	\par % каждый источник начинается с нового абзаца (будет стандартный отступ 1.25 см)
	% Записываем номер в aux-файл, чтобы команда \cite{} в тексте работала корректно
	\if@filesw \immediate\write\@auxout{\string\bibcite{#1}{\the\c@mybibcntr}}\fi
	\the\c@mybibcntr. \ignorespaces % выводим "1. " и игнорируем пробел после команды
}

\makeatother


%% Приложение ГОСТ
%% Счётчик приложений
\newcounter{gostappendix}
\renewcommand{\thegostappendix}{%
\ifcase\value{gostappendix}%
\or А\or Б\or В\or Г\or Д\or Е\or Ж\or И\or К%
\or Л\or М\or Н\or П\or Р\or С\or Т\or У\or Ф%
\or Х\or Ц\or Ч\or Ш\or Щ\or Э\or Ю\or Я%
\else ?\fi
}

%% Команда для приложения
\NewDocumentCommand{\gostappendix}{m}{%
\stepcounter{gostappendix}%
\clearpage
\phantomsection
	
% Сбрасываем счётчики для нового приложения
\setcounter{section}{0}%
\setcounter{figure}{0}%
\setcounter{table}{0}%
\setcounter{equation}{0}% % Формулы тоже по ГОСТу нумеруются в приложениях
	
% Задаём новый формат нумерации с буквой приложения
\renewcommand{\thesection}{\thegostappendix.\arabic{section}}%
\renewcommand{\thefigure}{\thegostappendix.\arabic{figure}}%
\renewcommand{\thetable}{\thegostappendix.\arabic{table}}%
\renewcommand{\theequation}{\thegostappendix.\arabic{equation}}%

\addcontentsline{toc}{section}{Приложение~\thegostappendix\ --- #1}%
\begin{center}
{\normalsize\bfseries ПРИЛОЖЕНИЕ~\thegostappendix\par}
{\bfseries #1\par}
\end{center}
}


%% Титульный лист
%% Определение заполнителей титульного листа
\newcommand{\UniversityDepartment}[1]{\gdef\setUniversityDepartment{#1}}
\newcommand{\setUniversityDepartment}{}
\newcommand{\TeacherPosition}[1]{\gdef\setTeacherPosition{#1}}
\newcommand{\setTeacherPosition}{}
\newcommand{\TeacherName}[1]{\gdef\setTeacherName{#1}}
\newcommand{\setTeacherName}{}
\newcommand{\WorkNumber}[1]{\gdef\setWorkNumber{#1}}
\newcommand{\setWorkNumber}{}
\newcommand{\WorkName}[1]{\gdef\setWorkName{#1}}
\newcommand{\setWorkName}{}
\newcommand{\UniversityDiscipline}[1]{\gdef\setUniversityDiscipline{#1}}
\newcommand{\setUniversityDiscipline}{}
\newcommand{\StudentGroup}[1]{\gdef\setStudentGroup{#1}}
\newcommand{\setStudentGroup}{}
\newcommand{\StudentName}[1]{\gdef\setStudentName{#1}}
\newcommand{\setStudentName}{}
\newcommand{\DATE}[1]{\gdef\setDATE{#1}}
\newcommand{\setDATE}{\today}
